\section{Details on Experiments}
\label{app:experiments}

\paragraph{Post-processing for Equalized Odds}
To derive classifiers that satisfy the Equalized Odds notion of
fairness, we use the method introduced by \citet{hardt2016equality}.
Essentially, the false-positive and false-negative constraints
are satisfied by randomly flipping some of the predictions of the
original classifiers. Let $q^{(t)}_\text{n2p}$ be the probability for group $G_t$ of
``flipping'' a negative prediction to positive,
and $q^{(t)}_\text{p2n}$ be that of flipping a positive prediction to negative.
The derived classifiers $\eo h_1$ and $\eo h_2$ essentially flip predictions according to
these probabilities:
%
\begin{align*}
  \eo h_t(\x) &= \begin{cases}
    \left( 1 - h_t(\x) \right) B^{(t)}_\text{p2n}
    + h_t(\x) \left(1 - B^{(t)}_\text{p2n} \right)
    & h_t(\x) \geq 0.5
    \\
    \left( 1 - h_t(\x) \right) B^{(t)}_\text{n2p}
    + h_t(\x) \left(1 - B^{(t)}_\text{n2p} \right)
    & h_t(\x) < 0.5
  \end{cases}
\end{align*}
%
where $B^{(t)}_\text{n2p}$ and $B^{(t)}_\text{p2n}$ are Bernoulli random variables
with expectations $q^{(t)}_\text{n2p}$ and $q^{(t)}_\text{p2n}$ respectively.
Note that this is a probabilistic generalization of the derived classifiers
presented in \cite{hardt2016equality}. If all outputs of $h_t$ were either
$0$ or $1$ we would arrive at the original formulation.

We can find the best rates $q^{(1)}_\text{n2p}$, $q^{(1)}_\text{p2n}$,
$q^{(2)}_\text{n2p}$, and $q^{(2)}_\text{p2n}$ through the following optimization
problem:
%
\begin{mini*}
{q^{(1)}_\text{}, q^{(1)}_\text{p2n}, q^{(2)}_\text{n2p}, q^{(2)}_\text{p2n}}
{\mathcal{L}(\eo h_1) + \mathcal{L}(\eo h_2)}
{}{}
\addConstraint{\fp{\eo h_1} = \fn{\eo h_1}}
\addConstraint{\fp{\eo h_2} = \fn{\eo h_2}}
\end{mini*}
%
where $\mathcal{L}$ represents the $0$/$1$ loss of the classifier:
%
\begin{align*}
  \mathcal L (h_t) &=
  \Pr_{G_t} \left[ h_t(\x) \geq 0.5 \mid y \!=\! 0 \right]
  +
  \Pr_{G_t} \left[ h_t(\x) < 0.5 \mid y \!=\! 1 \right].
\end{align*}
%
The two constraints
enforce the Equalized Odds constraints.
\citet{hardt2016equality} show that this can be solved via
a linear program.

\paragraph{Constrained-learning for Equalized Odds}
\citet{zafar2017fairness} introduce a method to achieve Equalized Odds
(under the name \emph{Disparate Mistreatment}) at training time using
optimization constraints. The problem is set up at learning a logistic
classifier under the Equalized Odds constraints. While these constraints
make the problem non-convex, \citeauthor{zafar2017fairness} show how
to formulate the problem as a disciplined convex-concave program.
Though this is generally intractable, it can be solved in many instances.
We refer the reader to \cite{zafar2017fairness} for details.

\paragraph{Training Procedure for Income Prediction.}
We train three models: a random forest, a multi-layer perceptron, and a SVM
with an RBF kernel. We convert the categorical features into one-hot enocodings.
$10\%$ of the data is reserved for hyperparameter tuning and
post-processing, and an additional $10\%$ is saved for final evaluation.
The random forest and MLP are naturally probabilistic and well calibrated.
We use Platt scaling to calibrate the SVM.
The hyperparameters were tuned by grid search based on 3-fold cross validation.
\autoref{fig:adult-res} displays the average false-positive and false-negative
costs across all models.

\paragraph{Training Procedure for Health Prediction.}
We train  a random forest and a linear SVM on this dataset. We use the
same dataset split and hyperparameter selection as with Income Prediction.
\autoref{fig:health-res} displays the average false-positive and false-negative
costs across all models.

\paragraph{Training Procedure for Health Prediction.}
For the trained Equalized Odds baseline we train a constrained
logistic classifier using the method proposed by \cite{zafar2017fairness}.
We derive the post-processed
classifiers (both for Equalized Odds and its calibrated relaxation) from the
original COMPAS classifier \cite{dieterich-northpointe-fairness}.
